\documentclass[10pt]{scrartcl}
\usepackage[mdthm,colorsec,noauthor,noasy]{evan}
\usepackage{handout}
\usepackage{figure-ink}
\newcommand{\vv}[1]{\mathbf{#1}}
\DeclareMathOperator{\spn}{Span}
\DeclareMathOperator{\Col}{Col}
\DeclareMathOperator{\Nul}{Nul}
\DeclareMathOperator{\tr}{tr}
% Method box: "How to ..." with numbered steps (standard, Teach every procedure).
\mdfdefinestyle{howtostyle}{nobreak=true,linewidth=0.8pt,linecolor=NavyBlue!70,backgroundcolor=NavyBlue!3,roundcorner=4pt,frametitlefont=\sffamily\bfseries\color{NavyBlue},frametitleaboveskip=5pt,frametitlebelowskip=1pt,innertopmargin=2pt,innerbottommargin=6pt,skipabove=9pt,skipbelow=4pt}
\newenvironment{howto}[1]{\begin{mdframed}[style=howtostyle,frametitle={#1}]}{\end{mdframed}}
\title{Math 54 --- Week 13}
\subtitle{Linear systems of differential equations: normal form, fundamental matrices, the eigenvalue method (Nagle 9.1, 9.4, 9.5)}
\author{}\date{}\ihead{}

\begin{document}
\maketitle

\section{Facts used from Weeks 1--12}
\begin{enumerate}[label=(W\arabic*),leftmargin=2.8em,itemsep=4pt]
  \ii \target{fact-W1}\textbf{Invertible matrices (\xlink{math54}{2026-09-14-matrix-algebra-inverses}{name-invertible-matrix-theorem}{Weeks~3--6}):} For square $M$: $\det M\neq 0\iff M$ invertible $\iff$ columns independent $\iff M\vv c=\vv 0$ only for $\vv c=\vv 0\iff M\vv c=\vv b$ has a unique solution for every $\vv b$. Scaling one column by $k$ multiplies $\det M$ by $k$, and $\det M^{\mathsf T}=\det M$ (\xlink{math54}{2026-09-21-determinants-vector-spaces}{thm-lay-3-5}{Week~5}). Triangular $\det$ equals the diagonal product (\xlink{math54}{2026-09-14-matrix-algebra-inverses}{thm-lay-3-2}{Week~4}). Rank Theorem: $\rank M+\dim\Nul M=n$ (\xlink{math54}{2026-09-28-bases-dimension-rank}{name-rank-theorem}{Week~6}).
  \ii \target{fact-W2}\textbf{Eigenvalues \& diagonalization (\xlink{math54}{2026-10-05-coordinates-eigenvalues}{thm-lay-5-2}{Weeks~7--8}):} $\lambda$ is an eigenvalue of $A\iff \det(A-\lambda I)=0$, with eigenvectors the nonzero vectors of $\Nul(A-\lambda I)$. Eigenvectors for distinct eigenvalues are linearly independent. $A=PDP^{-1}$ iff $A$ has $n$ linearly independent eigenvectors (\xlink{math54}{2026-10-12-diagonalization-complex-eigenvalues}{thm-lay-5-5}{Week~8}).
  \ii \target{fact-W3}\textbf{Symmetric matrices (\xlink{math54}{2026-11-02-symmetric-svd-second-order}{thm-lay-7-3}{Week~11}):} Any real symmetric matrix ($A^{\mathsf T}=A$) has real eigenvalues and $n$ orthonormal (hence independent) eigenvectors.
  \ii \target{fact-W4}\textbf{Scalar linear ODEs (\xlink{math54}{2026-11-02-symmetric-svd-second-order}{thm-nagle-4-2}{Weeks~11--12}):} Roots of $ar^2+br+c=0$ give $e^{rt}$; Wronskian of two solutions is never $0$ or always $0$; nonhomogeneous general solution is $y_p+y_h$ (\xlink{math54}{2026-11-09-undetermined-coefficients-superposition}{thm-nagle-4-4}{Week~12}). If a complex function solves a linear equation with real coefficients, so do its real and imaginary parts (\xlink{math54}{2026-11-02-symmetric-svd-second-order}{lemma-nagle-4-2}{Week~11, Real solutions from a complex solution}).
  \ii \target{fact-W5}\textbf{Calculus:} $\frac{d}{dt}e^{rt}=re^{rt}$; product/quotient/chain rules; for continuous $g$ the function $t\mapsto\int_{t_0}^t g(s)\,ds$ has derivative $g(t)$ (fundamental theorem of calculus); $f'\equiv 0$ on an interval $\implies f$ is constant; a continuous function on a closed bounded interval is bounded and attains its maximum (extreme value theorem); the absolute value of an integral is at most the integral of the absolute value; the series $\sum_k c^k/k!$ converges (to $e^c$), so its terms and its tails $\sum_{k\ge N}c^k/k!$ tend to $0$; vector/matrix derivatives and integrals act entry-by-entry.
\end{enumerate}

\section{This week}
\subsection{Systems in normal form (Nagle 9.1, 9.4)}
\begin{qbox}
  Two connected brine tanks hold salt amounts $x(t)$ and $y(t)$ satisfying
  \[
    x'=-4x+2y, \qquad y'=4x-4y.
  \]
  How do we write these two equations as one matrix equation? And how do we turn the second-order spring equation $my''+ky=0$ into a first-order system of the same kind?
\end{qbox}

The tank equations are already solved for the derivatives $x'$ and $y'$, and their right sides are linear combinations of $x$ and $y$. Each entry of $A\vv x$ is the dot product of a row of $A$ with $\vv x$, so the coefficients of $x$ and $y$ in the equation for $x'$ form row~1 of $A$, and those in the equation for $y'$ form row~2:
\[
  \begin{bmatrix}x'\\y'\end{bmatrix} = \begin{bmatrix}-4&2\\4&-4\end{bmatrix}\begin{bmatrix}x\\y\end{bmatrix}.
\]

\begin{fig}
\begin{tikzpicture}[scale=1.0,>=Stealth]
  % Coupled brine tanks: ink outlines, light blue brine stipple; every label black on a white backing.
  \InkRegion[inkblue]{(0.0,0.85) rectangle (1.85,2.20)}
  \InkStroke[black!85,line width=0.52pt]{(0.0,0.85) rectangle (1.85,2.20)}
  \InkRegion[inkblue]{(3.45,0.85) rectangle (5.30,2.20)}
  \InkStroke[black!85,line width=0.52pt]{(3.45,0.85) rectangle (5.30,2.20)}
  \node[font=\footnotesize,text=black,anchor=north] at (0.925,0.78) {Tank 1};
  \node[font=\footnotesize,text=black,anchor=north] at (4.375,0.78) {Tank 2};
  \node[font=\footnotesize,text=black,fill=white,inner sep=2pt,rounded corners=1pt] at (0.925,1.52) {salt $x(t)$};
  \node[font=\footnotesize,text=black,fill=white,inner sep=2pt,rounded corners=1pt] at (4.375,1.52) {salt $y(t)$};

  \draw[ink line,-{Stealth[length=3.8pt,width=2.2pt]},black,line width=0.55pt] (1.92,1.85) -- node[above=1pt,font=\footnotesize] {$4x$} (3.38,1.85);
  \draw[ink line,-{Stealth[length=3.8pt,width=2.2pt]},black,line width=0.55pt] (3.38,1.20) -- node[below=1pt,font=\footnotesize] {$2y$} (1.92,1.20);
  \draw[ink line,-{Stealth[length=3.8pt,width=2.2pt]},black,line width=0.55pt] (5.37,1.52) -- (6.15,1.52) node[right=2pt,font=\footnotesize] {drain $2y$};

  % Fresh-water inflow keeps the liquid volume constant and adds no salt.
  \draw[ink arrow,inkblue,line width=0.6pt] (.925,3.05)--(.925,2.26);
  \node[font=\footnotesize,anchor=west,text=black] at (1.08,2.70) {fresh water (no salt)};

  \node[font=\footnotesize,text=black] at (2.65,-0.05) {$x' = -4x + 2y,\qquad y' = 4x - 4y$};
\end{tikzpicture}
\figcap{The arrows give the rates at which salt moves; fresh water enters tank~1 and keeps both volumes constant. Rate in minus rate out for each tank gives the two equations under the tanks.}
\end{fig}

\noindent\textbf{Definition (Normal form, Nagle 9.4).} A system of $n$ linear differential equations is in \textbf{normal form} if
\[
  \vv x'(t) = A(t)\vv x(t) + \vv f(t),
\]
where $\vv x(t)=(x_1(t),\dots,x_n(t))^{\mathsf T}$, $A(t)$ is an $n\times n$ matrix function, and $\vv f(t)$ is a column vector function. It is \textbf{homogeneous} if $\vv f\equiv\vv 0$ and has \textbf{constant coefficients} if every entry of $A$ is constant. An \textbf{initial value problem (IVP)} specifies $\vv x(t_0)=\vv x_0$.

Normal form is the setting of the basic existence theorem. It is used below to prove that solutions are determined by their initial values and that $n$ independent solutions give all solutions. Its proof builds the solution by successive approximations; it is long, so it is given in \ilink{sec-proofs}{Section~3, Proofs of the facts used}.

\begin{theorem*}[Existence and uniqueness (Nagle 9.4, Thm.~2)]\target{thm-nagle-9-2}
  If $A(t)$ and $\vv f(t)$ are continuous on an open interval $I$ containing $t_0$, then for every $\vv x_0\in\RR^n$ the IVP $\vv x'=A(t)\vv x+\vv f(t)$, $\vv x(t_0)=\vv x_0$, has a unique solution on all of $I$.
\end{theorem*}
\noindent\emph{Proof.} See \ilink{proof-existence}{Proof of existence and uniqueness} in Section~3.

\begin{example*}[Reading off $A(t)$ when each equation is already solved for a derivative (Nagle 9.1)]\target{ex-read-off}
  Row $i$ of $A(t)$ lists the coefficients of $x_1,\dots,x_4$ in the equation for $x_i'$, with $0$ for a variable that does not appear:
  \[
    \left\{
    \begin{aligned}
      x_1'&=2x_1+t^2x_2+(4t+e^t)x_4,\\
      x_2'&=(\sin t)x_2+(\cos t)x_3,\\
      x_3'&=x_1+x_2+x_3+x_4,\quad x_4'=0
    \end{aligned}
    \right.
    \quad\Longleftrightarrow\quad
    \vv x'=
    \begin{bmatrix}
      2&t^2&0&4t+e^t\\0&\sin t&\cos t&0\\1&1&1&1\\0&0&0&0
    \end{bmatrix}\vv x.
  \]
\end{example*}

\medskip\noindent\textbf{Higher derivatives.} The spring equation contains $y''$, and normal form allows only first derivatives. The way around this is to give $y'$ its own name. With $x_1=y$ (position) and $x_2=y'$ (velocity), the definition $x_2=y'$ becomes the first-order equation $x_1'=x_2$, and $y''$ becomes $x_2'$, which the spring equation expresses through $x_1$. The same renaming works for any order:

\begin{proposition*}[Companion system of an $n$th-order equation (Nagle 9.1)]\target{prop-companion}
  Let $y^{(n)}+p_{n-1}(t)y^{(n-1)}+\dots+p_0(t)y=g(t)$. Put $x_1=y,\ x_2=y',\ \dots,\ x_n=y^{(n-1)}$. Then $y$ solves the equation if and only if $\vv x=(x_1,\dots,x_n)^{\mathsf T}$ solves $\vv x'=A(t)\vv x+\vv f(t)$ with
  \[
    A = \begin{bmatrix}0&1&0&\cdots&0\\0&0&1&\cdots&0\\\vdots&\vdots&\vdots&\ddots&\vdots\\0&0&0&\cdots&1\\-p_0&-p_1&-p_2&\cdots&-p_{n-1}\end{bmatrix},\qquad
    \vv f = \begin{bmatrix}0\\\vdots\\0\\g\end{bmatrix}.
  \]
\end{proposition*}
\begin{proof}
  \begin{steps}
    \ii \textbf{The first $n-1$ rows.} For $k<n$, $x_k'=(y^{(k-1)})'=y^{(k)}=x_{k+1}$. Row $k$ of $A$ therefore has a single $1$, in column $k+1$, and entry $k$ of $\vv f$ is $0$.
    \ii \textbf{The last row.} Solving the equation for its highest derivative gives $x_n'=y^{(n)}=-p_0x_1-p_1x_2-\dots-p_{n-1}x_n+g$, which is row $n$ of $A$ and entry $n$ of $\vv f$.
    \ii \textbf{Assume $\vv x$ solves the system.} The first $n-1$ rows say $x_2=x_1',\ x_3=x_2'=x_1'',\dots$, so $x_k=y^{(k-1)}$ with $y=x_1$. The last row then says exactly that $y$ solves the scalar equation. So scalar solutions $y$ and system solutions $\vv x=(y,y',\dots,y^{(n-1)})^{\mathsf T}$ correspond one-to-one. \qedhere
  \end{steps}
\end{proof}

\begin{fig}
\begin{tikzpicture}[>=Stealth]
  \node[font=\footnotesize,text=black] (ode) at (0,1.50)
    {$y^{(n)} + p_{n-1}y^{(n-1)} + \dots + p_0 y = g$};
  \node[font=\footnotesize,text=black] (sys) at (0,0)
    {$x_k' = x_{k+1}\ (k<n),\qquad x_n' = -\sum_{k=0}^{n-1} p_k x_{k+1} + g$};
  \draw[ink line,-{Stealth[length=3.8pt,width=2.2pt]},black,line width=0.55pt] (0,1.18) -- node[right=5pt,font=\footnotesize,text=black] {$\vv x = (y,\,y',\,\dots,\,y^{(n-1)})^{\mathsf T}$} (0,0.36);
\end{tikzpicture}
\figcap{Naming $y,y',\dots,y^{(n-1)}$ as the state variables turns one order-$n$ equation into $n$ first-order equations.}
\end{fig}

\noindent Applying the \ilink{thm-nagle-9-2}{existence and uniqueness theorem} to this companion system shows that a solution of $y^{(n)}+p_{n-1}y^{(n-1)}+\dots+p_0y=g$ (with continuous $p_k$, $g$) exists and is determined by $y(t_0),y'(t_0),\dots,y^{(n-1)}(t_0)$. This is the fact about $n$th-order equations used in Weeks~11--12.

The general procedure handles any mixture of orders and of equations that contain several highest derivatives at once.

\begin{howto}{How to put a system in normal form}\target{howto-normal-form}
  \begin{steps}
    \ii \textbf{Choose the state variables.} For each unknown function, find the highest derivative of it that appears anywhere in the system; call its order $m$. The state variables for that unknown are the function and its derivatives of orders $1,\dots,m-1$ (a first-order unknown contributes only itself). Number all of them $x_1,\dots,x_n$ in a fixed order. The size $n$ is the sum of the highest orders.
    \ii \textbf{Solve for the highest derivatives.} If an equation contains the highest derivative of only one unknown, solve it for that derivative, dividing by its coefficient. If several highest derivatives appear together in two or more equations, treat them as the unknowns of a linear algebraic system and solve that system by elimination; this needs its coefficient matrix to be invertible. If a leading coefficient $a(t)$ is $0$ at some $t$, divide only on an interval where $a(t)\neq 0$; the normal form holds on that interval only.
    \ii \textbf{Write the chain rows.} If the state variable $x_k$ is $y^{(j)}$ with $j<m-1$, then $x_k'=y^{(j+1)}$ is again a state variable $x_\ell$; write $x_k'=x_\ell$.
    \ii \textbf{Substitute and sort.} In each equation from step~2, replace every derivative of order below $m$ by its state variable. This gives $x_k'$ for the state variable $x_k=y^{(m-1)}$. In the equation for $x_i'$, the coefficient of $x_j$ goes into row $i$, column $j$ of $A(t)$, and every term that contains no unknown goes into entry $i$ of $\vv f(t)$. All other entries are $0$.
    \ii \textbf{Translate the initial conditions.} The values $y(t_0),y'(t_0),\dots,y^{(m-1)}(t_0)$ of each unknown are the entries of $\vv x(t_0)$, in the order of step~1. A normal-form IVP takes exactly one value per state variable.
  \end{steps}
  \textbf{Checks.} $A$ is $n\times n$ and its columns follow the order of the state variables. A derivative that is not the highest (such as $x'$ in $y'+x'+2y=0$) is a state variable and is only substituted. $\vv f$ holds only the terms without unknowns.
\end{howto}

\begin{example*}[One third-order equation with initial conditions]\target{ex-nf-third-order}
  Put $2y'''-4y'+6y=4\cos t$, $y(0)=1$, $y'(0)=0$, $y''(0)=-2$, in normal form.
  \begin{steps}
    \ii State variables: $x_1=y$, $x_2=y'$, $x_3=y''$ (the highest derivative is $y'''$, so $n=3$).
    \ii Divide by the leading coefficient $2$: $y'''=2y'-3y+2\cos t$.
    \ii Chain rows: $x_1'=x_2$, $x_2'=x_3$. Substituting: $x_3'=-3x_1+2x_2+2\cos t$.
    \ii Sort into $A$ and $\vv f$, and read off $\vv x(0)=(y(0),y'(0),y''(0))$:
    \[
      \vv x'=\begin{bmatrix}0&1&0\\0&0&1\\-3&2&0\end{bmatrix}\vv x+\begin{bmatrix}0\\0\\2\cos t\end{bmatrix},\qquad \vv x(0)=\begin{bmatrix}1\\0\\-2\end{bmatrix}.
    \]
  \end{steps}
\end{example*}

\begin{example*}[The spring $my''+ky=0$ as a first-order system (Nagle 9.1)]\target{ex-nf-spring}
  Putting $x_1=y$ (position) and $x_2=y'$ (velocity) gives $x_1'=x_2$ and, after dividing by $m$, $x_2'=y''=-\frac{k}{m}x_1$:
  \[
    \begin{bmatrix}x_1\\x_2\end{bmatrix}' = \begin{bmatrix}0&1\\-k/m&0\end{bmatrix}\begin{bmatrix}x_1\\x_2\end{bmatrix}.
  \]
\end{example*}

\begin{example*}[Two coupled springs: divide by the leading coefficient (Nagle 9.1)]\target{ex-nf-coupled-springs}
  For $2x''+6x-2y=0$, $y''+2y-2x=0$, each unknown has order $2$, so $n=4$ with $x_1=x$, $x_2=x'$, $x_3=y$, $x_4=y'$. Dividing the first equation by~$2$ gives $x''=-3x+y$, so $x_2'=-3x_1+x_3$; the second gives $x_4'=y''=2x_1-2x_3$:
  \[
    \begin{bmatrix}x_1\\x_2\\x_3\\x_4\end{bmatrix}' =
    \begin{bmatrix}0&1&0&0\\-3&0&1&0\\0&0&0&1\\2&0&-2&0\end{bmatrix}
    \begin{bmatrix}x_1\\x_2\\x_3\\x_4\end{bmatrix}.
  \]
\end{example*}

\begin{example*}[Mixed orders: $x''+3x'-y=\sin t$, $y'+x'+2y=0$]\target{ex-nf-mixed}
  Initial conditions: $x(0)=1$, $x'(0)=0$, $y(0)=2$.
  \begin{steps}
    \ii The highest derivatives are $x''$ (order $2$) and $y'$ (order $1$), so $n=2+1=3$ with $x_1=x$, $x_2=x'$, $x_3=y$.
    \ii Solve each equation for its highest derivative: $x''=-3x'+y+\sin t$ and $y'=-x'-2y$. The $x'$ in the second equation is not a highest derivative; it is the state variable $x_2$.
    \ii Chain row: $x_1'=x_2$. Substituting: $x_2'=-3x_2+x_3+\sin t$ and $x_3'=-x_2-2x_3$. Hence
    \[
      \vv x'=\begin{bmatrix}0&1&0\\0&-3&1\\0&-1&-2\end{bmatrix}\vv x+\begin{bmatrix}0\\\sin t\\0\end{bmatrix},\qquad \vv x(0)=\begin{bmatrix}1\\0\\2\end{bmatrix}.
    \]
    \ii The three given values are all the IVP needs. A value of $y'(0)$ cannot also be prescribed: the second equation fixes it as $y'(0)=-x'(0)-2y(0)=-4$.
  \end{steps}
\end{example*}

\begin{example*}[Coupled highest derivatives: $x''+y''=2x$, $x''-y''=4y$]\target{ex-nf-coupled-highest}
  \begin{steps}
    \ii Both equations contain both highest derivatives $x''$ and $y''$, so neither can be solved alone. Treat $x''$, $y''$ as unknowns of the algebraic system
    \[
      \begin{bmatrix}1&1\\1&-1\end{bmatrix}\begin{bmatrix}x''\\y''\end{bmatrix}=\begin{bmatrix}2x\\4y\end{bmatrix},
    \]
    whose coefficient matrix has determinant $-2\neq 0$. Adding the equations gives $2x''=2x+4y$; subtracting gives $2y''=2x-4y$. So $x''=x+2y$ and $y''=x-2y$.
    \ii With $x_1=x$, $x_2=x'$, $x_3=y$, $x_4=y'$: $x_1'=x_2$, $x_2'=x_1+2x_3$, $x_3'=x_4$, $x_4'=x_1-2x_3$, that is,
    \[
      \vv x'=\begin{bmatrix}0&1&0&0\\1&0&2&0\\0&0&0&1\\1&0&-2&0\end{bmatrix}\vv x.
    \]
  \end{steps}
\end{example*}

\begin{example*}[A leading coefficient that vanishes: $ty''+y'+y=0$]\target{ex-nf-vanishing}
  \begin{steps}
    \ii The leading coefficient $t$ is $0$ at $t=0$. On $(0,\infty)$ divide by $t$ to get $y''=-(y+y')/t$. With $x_1=y$, $x_2=y'$,
    \[
      \vv x'=\begin{bmatrix}0&1\\-1/t&-1/t\end{bmatrix}\vv x\qquad(t>0).
    \]
    The same matrix is the normal form on $(-\infty,0)$; the two intervals are separate problems.
    \ii $A(t)$ is continuous on $(0,\infty)$, so for $t_0=1$ and any $y(1)$, $y'(1)$ the \ilink{thm-nagle-9-2}{existence and uniqueness theorem} gives a unique solution on $(0,\infty)$. It says nothing about $t\le 0$.
    \ii At $t_0=0$ the IVP cannot be posed freely. Setting $t=0$ in the equation leaves $y'(0)+y(0)=0$, so for example $y(0)=1$, $y'(0)=1$ has no twice-differentiable solution.
  \end{steps}
\end{example*}

\medskip\noindent\textbf{Linearity of $L[\vv x]=\vv x'-A\vv x$.} For constants $a,b$,
\[
  L[a\vv x+b\vv y]=a\vv x'+b\vv y'-aA\vv x-bA\vv y=aL[\vv x]+bL[\vv y].
\]
Thus combinations of solutions of $L[\vv x]=\vv 0$ are again solutions, and \textbf{superposition} gives $L[c_1\vv x_1+c_2\vv x_2]=c_1\vv g_1+c_2\vv g_2$ when $L[\vv x_i]=\vv g_i$.

\revisit The tank equations are $\vv x'=A\vv x$ with $A=\begin{psmallmatrix}-4&2\\4&-4\end{psmallmatrix}$ and $\vv x=(x,y)^{\mathsf T}$: row~1 holds the coefficients of the equation for $x'$, row~2 those of the equation for $y'$. For $my''+ky=0$ the state variables are the position $x_1=y$ and the velocity $x_2=y'$, and the system is $\vv x'=\begin{psmallmatrix}0&1\\-k/m&0\end{psmallmatrix}\vv x$ (\ilink{ex-nf-spring}{Example: the spring}). An initial position $y(0)=y_0$ and velocity $y'(0)=v_0$ become $\vv x(0)=(y_0,v_0)^{\mathsf T}$.

\begin{shortcuts}
  \ii \textbf{Reading off $A$:} Row $i$ lists the coefficients of $x_1,\dots,x_n$ in the equation for $x_i'$.
  \ii \textbf{Companion matrix form:} $y''+by'+cy=0$ has $A=\begin{psmallmatrix}0&1\\-c&-b\end{psmallmatrix}$, whose characteristic polynomial $\det\begin{psmallmatrix}-r&1\\-c&-b-r\end{psmallmatrix}=r^2+br+c$ equals the auxiliary polynomial.
  \ii \textbf{System dimension:} Sum the highest derivative orders of all scalar unknowns ($2+2=4$ for the \ilink{ex-nf-coupled-springs}{coupled springs}, $2+1=3$ for the \ilink{ex-nf-mixed}{mixed-order example}).
  \ii \textbf{Coupled highest derivatives with a $\pm$ pattern:} add and subtract the equations, as in the \ilink{ex-nf-coupled-highest}{coupled highest derivatives example}.
\end{shortcuts}

\begin{trapbox}
  \ii \textbf{Divide out the leading coefficient first:} $2x''=-6x+2y$ gives $x_2'=-3x_1+x_3$, \emph{not} $-6x_1+2x_3$.
  \ii \textbf{Nonhomogeneous terms belong in $\vv f$, not $A$:} $x'=3x+t^2$ means $\vv x'=[3]\vv x+[t^2]$.
  \ii \textbf{Match column order to variable order:} With $(x,x',y,y')$, $y$ sits in column~$3$ (not~$2$).
  \ii \textbf{Too many initial conditions:} a first-order unknown gets one initial value. In the \ilink{ex-nf-mixed}{mixed-order example}, $y'(0)$ is determined by the equations and cannot be chosen.
  \ii \textbf{Highest derivatives that cannot be solved for:} in $x''+y''=x$, $x''+y''=y$ the coefficient matrix $\begin{psmallmatrix}1&1\\1&1\end{psmallmatrix}$ is singular. Subtracting gives $x=y$, so the system is really the single equation $2x''=x$ with $y=x$: its normal form has $2$ state variables, not $4$.
  \ii \textbf{A vanishing leading coefficient:} the normal form and the existence theorem hold only on an interval where that coefficient is nonzero; an IVP at a zero of it may have no solution (\ilink{ex-nf-vanishing}{$ty''+y'+y=0$}).
\end{trapbox}

\subsection{Independence, Wronskians, and fundamental matrices (Nagle 9.4)}
\begin{qbox}
  The three vector functions below solve $\vv x'=A\vv x$ for $A$ on the right. How can we prove they form a complete basis so that \emph{every} solution equals $c_1\vv x_1+c_2\vv x_2+c_3\vv x_3$?
  \[
    \vv x_1=e^{2t}\begin{bmatrix}1\\1\\1\end{bmatrix},\quad
    \vv x_2=e^{-t}\begin{bmatrix}-1\\0\\1\end{bmatrix},\quad
    \vv x_3=e^{-t}\begin{bmatrix}0\\1\\-1\end{bmatrix};\qquad
    A=\begin{bmatrix}0&1&1\\1&0&1\\1&1&0\end{bmatrix}.
  \]
\end{qbox}

\noindent\textbf{Definition (Linear dependence, Nagle 9.4 Def.~1).} Vector functions $\vv x_1,\dots,\vv x_m$ are \textbf{linearly dependent} on an interval $I$ if there exist \emph{constants} $c_1,\dots,c_m$, not all~$0$, such that $c_1\vv x_1(t)+\dots+c_m\vv x_m(t)=\vv 0$ for every $t\in I$; otherwise they are \textbf{linearly independent} on $I$.

\begin{example*}[A dependence relation with constant coefficients (Nagle 9.4)]\target{ex-dep-constant}
  $\vv x_1=(e^t,0,e^t)$, $\vv x_2=(3e^t,0,3e^t)$, $\vv x_3=(t,1,0)$ are dependent on $\RR$: the constants $3,-1,0$ give $3\vv x_1-\vv x_2+0\vv x_3\equiv\vv 0$.
\end{example*}

\begin{example*}[Independent although dependent at each single $t$ (Nagle 9.4)]\target{ex-abs-t}
  $\vv x_1=(t,t)$ and $\vv x_2=(|t|,|t|)$ are \textbf{independent} on $(-\infty,\infty)$ even though at every individual $t$ one is a scalar multiple of the other. Suppose $c_1\vv x_1+c_2\vv x_2\equiv\vv 0$ with constants $c_1,c_2$. At $t=1$ this says $c_1+c_2=0$, and at $t=-1$ it says $-c_1+c_2=0$, so $c_1=c_2=0$. The multiple that works at $t=1$ is not the one that works at $t=-1$, and the definition needs one set of constants for all $t$.
\end{example*}

\begin{example*}[Independence from one evaluation (Nagle 9.4)]\target{ex-indep-eval}
  $\vv x_1=e^{2t}(1,0,1)$, $\vv x_2=e^{2t}(1,1,-1)$, $\vv x_3=e^t(1,2,1)$ are independent on $\RR$. A relation $\sum c_i\vv x_i\equiv\vv 0$ holds in particular at $t=0$, which gives $M\vv c=\vv 0$ for
  \[
    M=\begin{bmatrix}1&1&1\\0&1&2\\1&-1&1\end{bmatrix},\qquad \det M=4\neq 0,
  \]
  so $c_1=c_2=c_3=0$ by (\ilink{fact-W1}{W1}).
\end{example*}

\noindent\textbf{Definition (Wronskian, Nagle 9.4 Def.~2).} The \textbf{Wronskian} of $n$ vector functions in $\RR^n$ is
\[
  W[\vv x_1,\dots,\vv x_n](t) = \det\bigl[\,\vv x_1(t)\ \cdots\ \vv x_n(t)\,\bigr].
\]
We write $X(t)=[\,\vv x_1(t)\ \cdots\ \vv x_n(t)\,]$, so $W(t)=\det X(t)$.

\begin{proposition*}[Wronskian test (Nagle 9.4)]\target{name-wronskian-test}
  If $W[\vv x_1,\dots,\vv x_n](t_0)\neq 0$ at even one point $t_0\in I$, then $\vv x_1,\dots,\vv x_n$ are linearly independent on $I$.
\end{proposition*}
\begin{proof}
  If $c_1\vv x_1(t)+\dots+c_n\vv x_n(t)\equiv\vv 0$ on $I$, setting $t=t_0$ gives $X(t_0)\vv c=\vv 0$ where $\det X(t_0)=W(t_0)\neq 0$; hence $\vv c=\vv 0$ by (\ilink{fact-W1}{W1}). \qedhere
\end{proof}

\Needspace{12\baselineskip}
The \ilink{ex-abs-t}{example with $(t,t)$ and $(|t|,|t|)$} has $W(t)=t|t|-|t|t\equiv 0$, yet the functions are independent: for arbitrary functions the converse of the Wronskian test fails. For solutions of \emph{one} linear system, the \ilink{thm-nagle-9-2}{existence and uniqueness theorem} makes $W$ either never zero or always zero:

\begin{theorem*}[Wronskian of solutions (Nagle 9.4)]\target{name-wronskian-solutions}
  Let $\vv x_1,\dots,\vv x_n$ solve $\vv x'=A(t)\vv x$ on $I$ with $A$ continuous. Either $W(t)\neq 0$ for every $t\in I$ (solutions independent on $I$), or $W(t)=0$ for every $t\in I$ (solutions dependent on $I$).
\end{theorem*}
\begin{ideabox}
  If $W(t_0)=0$ at one time $t_0$, the vectors $\vv x_i(t_0)$ are linearly dependent, so a nontrivial linear combination $\vv z=\sum c_i\vv x_i$ vanishes at $t_0$. Since $\vv z$ solves $\vv x'=A\vv x$ with $\vv z(t_0)=\vv 0$, and the zero function solves the same IVP, uniqueness forces $\vv z\equiv\vv 0$ at \emph{every} time $t$.
\end{ideabox}
\begin{proof}
  \begin{steps}
    \ii Suppose $W(t_0)=0$ at some $t_0\in I$. By (\ilink{fact-W1}{W1}), $X(t_0)\vv c=\vv 0$ for some $\vv c\neq\vv 0$.
    \ii Set $\vv z(t)=\sum c_i\vv x_i(t)$. By linearity of $L$, $\vv z'=A\vv z$ and $\vv z(t_0)=\vv 0$.
    \ii Since $\vv 0$ also solves this IVP, the \ilink{thm-nagle-9-2}{existence and uniqueness theorem} gives $\vv z(t)\equiv\vv 0$ on $I$: $\vv x_1,\dots,\vv x_n$ are dependent on $I$, and $X(t)\vv c=\vv 0$ with $\vv c\neq\vv 0$ gives $W(t)=\det X(t)=0$ for every $t\in I$.
    \ii If $W$ is nonzero at some point of $I$, step~3 shows it is nonzero at every point, and the \ilink{name-wronskian-test}{Wronskian test} gives independence. \qedhere
  \end{steps}
\end{proof}

\begin{howto}{How to decide whether vector functions are linearly independent on an interval}\target{howto-independence}
  \begin{steps}
    \ii \textbf{Count.} If there are $n$ functions with values in $\RR^n$, compute $W(t_0)$ at a convenient point $t_0$ (usually $0$). If $W(t_0)\neq 0$, they are independent (\ilink{name-wronskian-test}{Wronskian test}).
    \ii \textbf{If $W(t_0)=0$.} If the functions are known to solve one system $\vv x'=A(t)\vv x$ with $A$ continuous on $I$, they are dependent (\ilink{name-wronskian-solutions}{Wronskian of solutions}). Otherwise try other points: a single nonzero value settles independence.
    \ii \textbf{If $W\equiv 0$ and nothing is known about a system,} look for constants. Evaluate $\sum c_i\vv x_i(t)=\vv 0$ at several values of $t$; if the resulting equations force $\vv c=\vv 0$, the functions are independent (\ilink{ex-abs-t}{as for $(t,t)$, $(|t|,|t|)$}). If you find constants that work for every $t$, they are dependent (\ilink{ex-dep-constant}{as in this example}).
    \ii \textbf{Fewer than $n$ functions} have no Wronskian. Use step~3: evaluate at one or more points and solve for the $c_i$.
  \end{steps}
\end{howto}

\begin{example*}[Two vector functions in $\RR^3$: no Wronskian]\target{ex-two-in-r3}
  For $\vv x_1=(1,t,t^2)$ and $\vv x_2=(t,1,0)$, the relation $c_1\vv x_1+c_2\vv x_2\equiv\vv 0$ at $t=0$ reads $c_1(1,0,0)+c_2(0,1,0)=\vv 0$, so $c_1=c_2=0$: the functions are independent on $\RR$.
\end{example*}

\begin{theorem*}[Representation, homogeneous case (Nagle 9.4, Thm.~3)]\target{thm-nagle-9-3}
  If $A(t)$ is continuous on $I$ and $\vv x_1,\dots,\vv x_n$ are linearly independent solutions of $\vv x'=A(t)\vv x$ on $I$, then every solution on $I$ is $\vv x=c_1\vv x_1+\dots+c_n\vv x_n=X(t)\vv c$ for unique constants $c_1,\dots,c_n$.
\end{theorem*}
\begin{proof}
  \begin{steps}
    \ii Let $\vv x$ be any solution and fix $t_0\in I$. By \ilink{name-wronskian-solutions}{Wronskian of solutions}, $\det X(t_0)=W(t_0)\neq 0$, so $X(t_0)\vv c=\vv x(t_0)$ has a unique solution $\vv c$ by (\ilink{fact-W1}{W1}).
    \ii The function $\vv y(t)=X(t)\vv c$ solves $\vv y'=A\vv y$ (linearity of $L$) and $\vv y(t_0)=\vv x(t_0)$. The \ilink{thm-nagle-9-2}{existence and uniqueness theorem} forces $\vv x\equiv\vv y=X(t)\vv c$. \qedhere
  \end{steps}
\end{proof}

Such $\vv x_1,\dots,\vv x_n$ form a \textbf{fundamental solution set}, $X(t)=[\,\vv x_1(t)\ \cdots\ \vv x_n(t)\,]$ is a \textbf{fundamental matrix} (invertible for all $t\in I$ and obeying $X'=AX$, because column $i$ of $AX$ is $A\vv x_i=\vv x_i'$), and $\vv x=X(t)\vv c$ is the \textbf{general solution}. Fundamental sets always exist: for fixed $t_0\in I$, let $\vv x_i$ solve the IVP $\vv x_i(t_0)=\vv e_i$, which the \ilink{thm-nagle-9-2}{existence and uniqueness theorem} allows; then $W(t_0)=\det I_n=1\neq 0$.

\begin{theorem*}[Representation, nonhomogeneous case (Nagle 9.4, Thm.~4)]\target{thm-nagle-9-4}
  If $A(t)$ and $\vv f(t)$ are continuous on $I$, $\vv x_p$ solves $\vv x'=A\vv x+\vv f$ on $I$, and $X(t)$ is a fundamental matrix for $\vv x'=A\vv x$, then every solution on $I$ has the form $\vv x=\vv x_p+X(t)\vv c$.
\end{theorem*}
\begin{proof}
  \begin{steps}
    \ii \textbf{Every solution has this form.} For any solution $\vv x$, $L[\vv x-\vv x_p]=L[\vv x]-L[\vv x_p]=\vv f-\vv f=\vv 0$, so $\vv x-\vv x_p=X\vv c$ by \ilink{thm-nagle-9-3}{Thm.~3}.
    \ii \textbf{Every function of this form is a solution.} $L[\vv x_p+X\vv c]=L[\vv x_p]+L[X\vv c]=\vv f+\vv 0=\vv f$. \qedhere
  \end{steps}
\end{proof}

\begin{howto}{How to check a fundamental set and solve an IVP with it}\target{howto-fundamental-ivp}
  \begin{steps}
    \ii \textbf{Each function is a solution.} Differentiate $\vv x_i$ and compare with $A\vv x_i$. For a nonhomogeneous system also check $\vv x_p'=A\vv x_p+\vv f$.
    \ii \textbf{Independence.} Compute $W(t_0)=\det X(t_0)$ at one convenient point; nonzero means fundamental.
    \ii \textbf{General solution.} $\vv x=X(t)\vv c$, or $\vv x=\vv x_p+X(t)\vv c$ for a nonhomogeneous system (\ilink{thm-nagle-9-4}{Thm.~4}).
    \ii \textbf{Initial condition $\vv x(t_0)=\vv x_0$.} Solve $X(t_0)\vv c=\vv x_0-\vv x_p(t_0)$ at the given $t_0$ (not at $0$ unless $t_0=0$), and substitute $\vv c$.
  \end{steps}
\end{howto}

\begin{example*}[A nonhomogeneous IVP from a given fundamental set]\target{ex-nonhom-ivp}
  Solve $\vv x'=A\vv x+\vv f$, $\vv x(0)=\vv 0$, for
  \[
    A=\begin{bmatrix}2&-1\\3&-2\end{bmatrix},\qquad \vv f=\begin{bmatrix}0\\1\end{bmatrix},\qquad
    \vv x_1=e^{t}\begin{bmatrix}1\\1\end{bmatrix},\quad \vv x_2=e^{-t}\begin{bmatrix}1\\3\end{bmatrix},\quad \vv x_p=\begin{bmatrix}1\\2\end{bmatrix}.
  \]
  \begin{steps}
    \ii Check: $A(1,1)^{\mathsf T}=(1,1)^{\mathsf T}$, so $\vv x_1'=\vv x_1=A\vv x_1$. $A(1,3)^{\mathsf T}=(-1,-3)^{\mathsf T}$, so $\vv x_2'=-\vv x_2=A\vv x_2$. $A\vv x_p=(0,-1)^{\mathsf T}$, so $A\vv x_p+\vv f=\vv 0=\vv x_p'$.
    \ii $W(0)=\det\begin{psmallmatrix}1&1\\1&3\end{psmallmatrix}=2\neq 0$, so $\vv x=\vv x_p+c_1\vv x_1+c_2\vv x_2$ is the general solution.
    \ii At $t=0$: $c_1+c_2=-1$ and $c_1+3c_2=-2$, so $c_2=-\tfrac12$ and $c_1=-\tfrac12$:
    \[
      \vv x(t)=\begin{bmatrix}1\\2\end{bmatrix}-\frac12e^{t}\begin{bmatrix}1\\1\end{bmatrix}-\frac12e^{-t}\begin{bmatrix}1\\3\end{bmatrix}.
    \]
  \end{steps}
\end{example*}

\begin{proposition*}[Abel's formula (Nagle 9.4, Prob.~32)]\target{prop-abel-formula}
  If $\vv x_1,\vv x_2$ solve $\vv x'=A(t)\vv x$ with continuous $2\times 2$ $A(t)$, then $W'=(\tr A)W$, so
  \[
    W(t) = W(t_0)\exp\left(\int_{t_0}^t \tr A(s)\,ds\right).
  \]
  The same formula, with $\operatorname{tr}A=a_{11}+\dots+a_{nn}$, holds for $n\times n$ systems (\ilink{proof-abel-n}{proof for $n\times n$} in Section~3).
\end{proposition*}
\begin{proof}[Proof for $2\times 2$ systems]
  \begin{steps}
    \ii Write $\vv x_1=(p,q)^{\mathsf T}$, $\vv x_2=(u,v)^{\mathsf T}$, $A=\begin{psmallmatrix}a&b\\c&d\end{psmallmatrix}$, so $W=pv-qu$ and $p'=ap+bq$, $q'=cp+dq$, $u'=au+bv$, $v'=cu+dv$. Substituting,
    \[
      W' = p'v+pv'-q'u-qu' = (a+d)(pv-qu) = (\tr A)W.
    \]
    \ii Put $E(t)=\exp\int_{t_0}^t\tr A(s)\,ds$. Then $E'=(\tr A)E$, $E(t_0)=1$, and $(W/E)'=(W'E-WE')/E^2=0$, so $W(t)=W(t_0)E(t)$. Since $E(t)>0$, this shows again that $W$ is either never $0$ or identically $0$. \qedhere
  \end{steps}
\end{proof}

\noindent\textbf{Link with the scalar Wronskian:} For $y''+by'+cy=0$, the companion matrix has $\tr A=-b$, so Abel's formula reads $W(t)=W(t_0)e^{-b(t-t_0)}$.

\revisit Each function solves the system: for example $A(1,1,1)^{\mathsf T}=(2,2,2)^{\mathsf T}$, so $A\vv x_1=2\vv x_1=\vv x_1'$, and $A(-1,0,1)^{\mathsf T}=(1,0,-1)^{\mathsf T}=-(-1,0,1)^{\mathsf T}$, $A(0,1,-1)^{\mathsf T}=(0,-1,1)^{\mathsf T}$, so $A\vv x_2=-\vv x_2=\vv x_2'$ and $A\vv x_3=-\vv x_3=\vv x_3'$. At $t=0$,
\[
  W(0) = \det\begin{bmatrix}1&-1&0\\1&0&1\\1&1&-1\end{bmatrix} = -3\neq 0,
\]
so $\vv x_1,\vv x_2,\vv x_3$ are independent by the \ilink{name-wronskian-test}{Wronskian test}, hence a fundamental set, and by \ilink{thm-nagle-9-3}{Thm.~3} every solution is $\vv x=c_1\vv x_1+c_2\vv x_2+c_3\vv x_3$. (Directly, $W(t)=e^{2t}e^{-t}e^{-t}W(0)=-3$, as \ilink{prop-abel-formula}{Abel's formula} predicts from $\tr A=0$.)

\begin{shortcuts}
  \ii \textbf{Test independence at $t=0$:} for solutions of one system, one determinant $W(0)\neq 0$ settles everything (\ilink{name-wronskian-test}{Wronskian test} and \ilink{name-wronskian-solutions}{Wronskian of solutions}).
  \ii \textbf{Exponential columns determinant:} $\det[\,e^{r_1 t}\vv u_1\ \cdots\ e^{r_n t}\vv u_n\,]=e^{(r_1+\dots+r_n)t}\det[\,\vv u_1\ \cdots\ \vv u_n\,]$ by (\ilink{fact-W1}{W1}).
  \ii \textbf{Abel's formula for constant $A$:} $W(t)=W(0)e^{(\tr A)t}$.
  \ii \textbf{IVP formula via fundamental matrix:} $\vv x(t)=X(t)X(t_0)^{-1}\vv x_0$ solves $\vv x'=A\vv x,\ \vv x(t_0)=\vv x_0$.
\end{shortcuts}

\begin{trapbox}
  \ii \textbf{Dependence at each $t$ is not dependence on the interval:} $(t,t)$ and $(|t|,|t|)$ are independent on $\RR$ despite $W\equiv 0$.
  \ii \textbf{$W\equiv 0$ proves dependence \emph{only} for solutions of \emph{one} continuous system:} \ilink{name-wronskian-solutions}{Wronskian of solutions} does not apply to functions like $(1,0),(0,t)$ ($W(0)=0$, independent) or $(t,1),(t^2,t)$ ($W\equiv 0$, independent): neither pair solves a system $\vv x'=A(t)\vv x$ with $A$ continuous across $t=0$.
  \ii \textbf{Nonhomogeneous IVPs/general solutions use $\vv x_p+X\vv c$:} the sum of two solutions of $\vv x'=A\vv x+\vv f$ satisfies $\vv x'=A\vv x+2\vv f$, not $\vv x'=A\vv x+\vv f$.
  \ii \textbf{Subtract $\vv x_p(t_0)$ before solving for $\vv c$:} the IVP equation is $X(t_0)\vv c=\vv x_0-\vv x_p(t_0)$.
\end{trapbox}

\subsection{Constant coefficients: the eigenvalue method (Nagle 9.5)}
\begin{qbox}
  Find all solutions of $\vv x'=A\vv x$ for
  \[
    A=\begin{bmatrix}2&-3\\1&-2\end{bmatrix}.
  \]
  Which solutions move along straight lines through the origin, and how does every trajectory behave as $t\to\infty$?
\end{qbox}

For scalar equations $e^{rt}$ worked; try the vector form $\vv x(t)=e^{rt}\vv u$ with constant $\vv u\neq\vv 0$. Since $\vv x'=re^{rt}\vv u$ and $e^{rt}\neq 0$:
\[
  \vv x'=A\vv x \iff A\vv u=r\vv u \iff (A-rI)\vv u=\vv 0.
\]
So $e^{rt}\vv u$ is a solution iff $r$ is an eigenvalue of $A$ (root of the \textbf{characteristic equation} $\det(A-rI)=0$) and $\vv u$ is an eigenvector for $r$ (\ilink{fact-W2}{W2}; Nagle 9.5 Def.~3).

\begin{example*}[Eigenpairs of the $2\times 2$ matrix in the question (Nagle 9.5)]\target{ex-eig-2x2}
  \[
    \det(A-rI)=(2-r)(-2-r)+3=r^2-1=0\implies r=1,\ -1.
  \]
  For $r=1$: $A-I$ has rows $(1,-3),(1,-3)$, so $u_1=3u_2$ and $\vv u_1=(3,1)$. For $r=-1$: $A+I$ has rows $(3,-3),(1,-1)$, so $u_1=u_2$ and $\vv u_2=(1,1)$.
\end{example*}

\begin{example*}[Eigenpairs of a $3\times 3$ matrix with distinct eigenvalues (Nagle 9.5)]\target{ex-eig-3x3}
  \[
    A = \begin{bmatrix}1&2&-1\\1&0&1\\4&-4&5\end{bmatrix}
    \implies \det(A-rI)=-(r-1)(r-2)(r-3).
  \]
  Row-reducing $A-rI$ gives eigenvectors $(-1,1,2)$ for $r=1$, $(-2,1,4)$ for $r=2$, and $(-1,1,4)$ for $r=3$ (check: $A(-2,1,4)=2(-2,1,4)$, $A(-1,1,4)=3(-1,1,4)$).
\end{example*}

\begin{theorem*}[$n$ independent eigenvectors (Nagle 9.5, Thm.~5)]\target{thm-nagle-9-5}
  If $n\times n$ matrix $A$ has $n$ linearly independent eigenvectors $\vv u_1,\dots,\vv u_n$ with eigenvalues $r_1,\dots,r_n$ (not necessarily distinct), then $\{e^{r_1 t}\vv u_1,\dots,e^{r_n t}\vv u_n\}$ is a fundamental solution set on $(-\infty,\infty)$, $X(t)=[\,e^{r_1 t}\vv u_1\ \cdots\ e^{r_n t}\vv u_n\,]$ is a fundamental matrix, and the general solution is
  \[
    \vv x(t) = c_1 e^{r_1 t}\vv u_1 + \dots + c_n e^{r_n t}\vv u_n.
  \]
\end{theorem*}
\begin{proof}
  Each $e^{r_i t}\vv u_i$ solves $\vv x'=A\vv x$ by the computation above. Factoring $e^{r_i t}$ out of column~$i$ (\ilink{fact-W1}{W1}) gives $W(t)=e^{(r_1+\dots+r_n)t}\det[\,\vv u_1\ \cdots\ \vv u_n\,]$, which is nonzero because $\vv u_1,\dots,\vv u_n$ are independent. So the solutions are independent by the \ilink{name-wronskian-test}{Wronskian test}, and every solution is a combination of them by \ilink{thm-nagle-9-3}{Thm.~3}. \qedhere
\end{proof}
\noindent \ilink{thm-nagle-9-5}{Thm.~5} also holds when some $r_i$ and $\vv u_i$ are complex; then $e^{r_i t}\vv u_i$ is a complex-valued solution. Its real and imaginary parts are real solutions, as for scalar equations (\ilink{fact-W4}{W4}): if $\vv z=\vv a+i\vv b$ solves $\vv z'=A\vv z$ with $A$ real, then $\vv a'+i\vv b'=A\vv a+iA\vv b$, and comparing real and imaginary parts gives $\vv a'=A\vv a$ and $\vv b'=A\vv b$.

\begin{theorem*}[Distinct eigenvalues guarantee independence (Nagle 9.5, Thm.~6, Cor.~1)]\target{thm-nagle-9-6}
  Eigenvectors $\vv u_1,\dots,\vv u_m$ corresponding to distinct eigenvalues $r_1,\dots,r_m$ are linearly independent. Hence any $n\times n$ matrix with $n$ distinct eigenvalues has fundamental set $\{e^{r_1 t}\vv u_1,\dots,e^{r_n t}\vv u_n\}$.
\end{theorem*}
\begin{ideabox}
  Two eigenvectors show the pattern. If $\vv u_1=c\vv u_2$, multiplying by $A$ gives $r_1\vv u_1=cr_2\vv u_2=r_2\vv u_1$, so $(r_1-r_2)\vv u_1=\vv 0$, which is impossible because $\vv u_1\neq\vv 0$ and $r_1\neq r_2$. Applying $A$ scales each eigenvector by its own eigenvalue, so a relation among eigenvectors turns into a second relation with different weights; subtracting the right multiple removes one vector. For $m$ vectors, use $A-r_{p+1}I$ to remove $\vv u_{p+1}$ from the first relation that appears.
\end{ideabox}
\begin{proof}
  \begin{steps}
    \ii Suppose the list is dependent. Since $\vv u_1\neq\vv 0$, there is a first index $p+1$ for which $\vv u_{p+1}$ is a combination of the earlier vectors: $\vv u_{p+1}=\sum_{i=1}^p c_i\vv u_i$, with $\vv u_1,\dots,\vv u_p$ independent.
    \ii Apply $A-r_{p+1}I$. The left side becomes $\vv 0$ because $A\vv u_{p+1}=r_{p+1}\vv u_{p+1}$. Each $\vv u_i$ on the right becomes $(r_i-r_{p+1})\vv u_i$. So $\sum_{i=1}^p c_i(r_i-r_{p+1})\vv u_i=\vv 0$.
    \ii Since $\vv u_1,\dots,\vv u_p$ are independent, $c_i(r_i-r_{p+1})=0$ for each $i$, and $r_i\neq r_{p+1}$ gives $c_i=0$. Then $\vv u_{p+1}=\vv 0$, which contradicts that eigenvectors are nonzero.
    \ii With $n$ distinct eigenvalues we get $n$ independent eigenvectors, and \ilink{thm-nagle-9-5}{Thm.~5} applies. \qedhere
  \end{steps}
\end{proof}

\revisit By the \ilink{ex-eig-2x2}{eigenpairs computed above} and \ilink{thm-nagle-9-6}{Thm.~6}, the general solution of $\vv x'=\begin{psmallmatrix}2&-3\\1&-2\end{psmallmatrix}\vv x$ is
\[
  \vv x(t) = c_1 e^t\begin{bmatrix}3\\1\end{bmatrix} + c_2 e^{-t}\begin{bmatrix}1\\1\end{bmatrix}.
\]
The straight-line trajectories are the single-term solutions $c_1e^t(3,1)$, which move outward from $\vv 0$ along the line through $(3,1)$, and $c_2e^{-t}(1,1)$, which decay to $\vv 0$ along the line through $(1,1)$. For every initial state with $c_1\neq 0$, the term $e^t$ dominates as $t\to\infty$, so the trajectory grows without bound and its direction approaches $\pm(3,1)$. Only the solutions with $c_1=0$ tend to $\vv 0$.

\begin{fig}
\begin{tikzpicture}[
  scale=1.0,>=Stealth,
  traj/.style={
    black!72, line width=0.46pt,
    postaction={decorate},
    decoration={markings, mark=at position #1 with {\arrow[black!85]{Stealth[length=3.3pt,width=2.0pt]}}}
  }
]
  % (a) Saddle of the question
  \begin{scope}[xshift=2.85cm,yshift=2.25cm,scale=0.68]
    \draw[ink line,figaxis] (-3.50,0) -- (3.50,0) node[right=1pt,font=\scriptsize] {\InkLabel{$x_1$}};
    \draw[ink line,figaxis] (0,-2.40) -- (0,2.40) node[above=1pt,font=\scriptsize] {\InkLabel{$x_2$}};
    % Exact trajectories x(t) = c1 e^t (3,1) + c2 e^-t (1,1) clipped to box
    \begin{scope}
      \clip (-3.25,-2.10) rectangle (3.25,2.10);
      \foreach \s in {1,-1} {
        \draw[traj=0.52] plot[domain=-1.55:1.65,samples=40] ({\s*(0.20*exp(\x)*3 + 0.48*exp(-\x)*1)}, {\s*(0.20*exp(\x)*1 + 0.48*exp(-\x)*1)});
      }
      \foreach \s in {1,-1} {
        \draw[traj=0.52] plot[domain=-1.45:1.72,samples=40] ({\s*(0.20*exp(\x)*3 - 0.55*exp(-\x)*1)}, {\s*(0.20*exp(\x)*1 - 0.55*exp(-\x)*1)});
        \draw[traj=0.52] plot[domain=-0.72:1.05,samples=35] ({\s*(0.45*exp(\x)*3 - 1.25*exp(-\x)*1)}, {\s*(0.45*exp(\x)*1 - 1.25*exp(-\x)*1)});
      }
    \end{scope}
    % Eigendirections: red for the growing direction u1, black for the decaying direction u2
    \InkStroke[figred,line width=0.75pt,-{Stealth[length=4.0pt,width=2.3pt]}]{(0,0) -- (3.25,1.083)}
    \InkStroke[figred,line width=0.75pt,-{Stealth[length=4.0pt,width=2.3pt]}]{(0,0) -- (-3.25,-1.083)}
    \InkStroke[black,line width=0.75pt,-{Stealth[length=4.0pt,width=2.3pt]}]{(2.10,2.10) -- (0.16,0.16)}
    \InkStroke[black,line width=0.75pt,-{Stealth[length=4.0pt,width=2.3pt]}]{(-2.10,-2.10) -- (-0.16,-0.16)}
    \fill[black] (0,0) circle (1.6pt);
    \node[font=\scriptsize,text=black] at (0.40,-0.30) {\InkLabel{$\vv 0$}};
    \node[font=\scriptsize,text=black,anchor=south west] at (2.05,2.08) {\InkLabel{$\vv u_2\;(r_2<0)$}};
    \node[font=\scriptsize,text=figred,anchor=west] at (3.28,1.08) {\InkLabel{$\vv u_1\;(r_1>0)$}};
  \end{scope}

  % (b) Sink (r2 < r1 < 0): all black
  \begin{scope}[xshift=8.85cm,yshift=2.25cm,scale=0.68]
    \draw[ink line,figaxis] (-2.25,0) -- (2.25,0) node[right=1pt,font=\scriptsize] {\InkLabel{$x_1$}};
    \draw[ink line,figaxis] (0,-2.35) -- (0,2.35) node[above=1pt,font=\scriptsize] {\InkLabel{$x_2$}};
    \begin{scope}
      \clip (-2.00,-1.88) rectangle (2.00,1.88);
      \foreach \sx/\sy in {1/1, -1/1, 1/-1, -1/-1} {
        \draw[traj=0.50] plot[domain=1.30:0.06,samples=30] ({\sx*\x}, {\sy*0.55*\x*\x*\x});
        \draw[traj=0.50] plot[domain=0.93:0.06,samples=30] ({\sx*\x}, {\sy*2.40*\x*\x*\x});
      }
    \end{scope}
    \InkStroke[black,line width=0.65pt,-{Stealth[length=4.0pt,width=2.3pt]}]{(2.00,0) -- (0.16,0)}
    \InkStroke[black,line width=0.65pt,-{Stealth[length=4.0pt,width=2.3pt]}]{(-2.00,0) -- (-0.16,0)}
    \InkStroke[black,line width=0.65pt,-{Stealth[length=4.0pt,width=2.3pt]}]{(0,1.95) -- (0,0.16)}
    \InkStroke[black,line width=0.65pt,-{Stealth[length=4.0pt,width=2.3pt]}]{(0,-1.95) -- (0,-0.16)}
    \fill[black] (0,0) circle (1.6pt);
    \node[font=\scriptsize,text=black,anchor=south] at (-1.82,0.06) {\InkLabel{$r_1$ (slow)}};
    \node[font=\scriptsize,text=black,anchor=west] at (0.10,2.12) {\InkLabel{$r_2$ (fast)}};
  \end{scope}

  % (c) Source (0 < r1 < r2): red on the faster direction r2
  \begin{scope}[xshift=13.75cm,yshift=2.25cm,scale=0.68]
    \draw[ink line,figaxis] (-2.25,0) -- (2.25,0) node[right=1pt,font=\scriptsize] {\InkLabel{$x_1$}};
    \draw[ink line,figaxis] (0,-2.35) -- (0,2.35) node[above=1pt,font=\scriptsize] {\InkLabel{$x_2$}};
    \begin{scope}
      \clip (-2.00,-1.88) rectangle (2.00,1.88);
      \foreach \sx/\sy in {1/1, -1/1, 1/-1, -1/-1} {
        \draw[traj=0.52] plot[domain=0.06:1.30,samples=30] ({\sx*\x}, {\sy*0.55*\x*\x*\x});
        \draw[traj=0.52] plot[domain=0.06:0.93,samples=30] ({\sx*\x}, {\sy*2.40*\x*\x*\x});
      }
    \end{scope}
    \InkStroke[black,line width=0.65pt,-{Stealth[length=4.0pt,width=2.3pt]}]{(0,0) -- (2.00,0)}
    \InkStroke[black,line width=0.65pt,-{Stealth[length=4.0pt,width=2.3pt]}]{(0,0) -- (-2.00,0)}
    \InkStroke[figred,line width=0.75pt,-{Stealth[length=4.0pt,width=2.3pt]}]{(0,0) -- (0,1.95)}
    \InkStroke[figred,line width=0.75pt,-{Stealth[length=4.0pt,width=2.3pt]}]{(0,0) -- (0,-1.95)}
    \fill[black] (0,0) circle (1.6pt);
    \node[font=\scriptsize,text=black,anchor=south] at (-1.82,0.06) {\InkLabel{$r_1$ (slow)}};
    \node[font=\scriptsize,text=figred,anchor=west] at (0.10,2.12) {\InkLabel{$r_2$ (fast)}};
  \end{scope}

  % Panel titles and plain roman panel labels
  \node[font=\scriptsize,anchor=south] at (2.85,4.28) {\InkLabel{Saddle ($r_1>0>r_2$)}};
  \node[font=\scriptsize,anchor=south] at (8.85,4.28) {\InkLabel{Sink ($r_2<r_1<0$)}};
  \node[font=\scriptsize,anchor=south] at (13.75,4.28) {\InkLabel{Source ($0<r_1<r_2$)}};
  \node[font=\footnotesize,anchor=north] at (2.85,0.22) {\InkLabel{(a)}};
  \node[font=\footnotesize,anchor=north] at (8.85,0.22) {\InkLabel{(b)}};
  \node[font=\footnotesize,anchor=north] at (13.75,0.22) {\InkLabel{(c)}};
\end{tikzpicture}
\figcap{Phase portraits by eigenvalue signs. \textbf{(a)} The saddle of the question: trajectories come in along $\vv u_2$ (black) and leave along $\vv u_1$ (red). \textbf{(b)} A sink approaches $\vv 0$ tangent to the slower eigendirection; \textbf{(c)} a source leaves $\vv 0$ tangent to the slower direction and turns toward the faster one (red).}
\end{fig}

\begin{example*}[A $3\times 3$ IVP (Nagle 9.5)]\target{ex-ivp-3x3}
  Solve $\vv x'=A\vv x$, $\vv x(0)=(-1,0,0)$, for the matrix of the \ilink{ex-eig-3x3}{$3\times 3$ eigenpair example}. The general solution is $c_1e^t(-1,1,2)+c_2e^{2t}(-2,1,4)+c_3e^{3t}(-1,1,4)$. At $t=0$, solving $c_1(-1,1,2)+c_2(-2,1,4)+c_3(-1,1,4)=(-1,0,0)$ gives $c_1=0$, $c_2=1$, $c_3=-1$, so
  \[
    \vv x(t)=e^{2t}\begin{bmatrix}-2\\1\\4\end{bmatrix}-e^{3t}\begin{bmatrix}-1\\1\\4\end{bmatrix}.
  \]
\end{example*}

\noindent\textbf{Symmetric matrices (Nagle 9.5 Def.~4).} Every symmetric $A$ ($A^{\mathsf T}=A$) has $n$ independent eigenvectors even when eigenvalues repeat (\ilink{fact-W3}{W3}), so \ilink{thm-nagle-9-5}{Thm.~5} always applies.

\begin{example*}[A symmetric matrix with a repeated eigenvalue (Nagle 9.5)]\target{ex-sym-repeated}
  \[
    A = \begin{bmatrix}1&-2&2\\-2&1&2\\2&2&1\end{bmatrix}, \qquad \det(A-rI)=-(r-3)^2(r+3).
  \]
  For $r=3$, $A-3I$ has rank~$1$ (every row is a multiple of $(-1,-1,1)$, so the eigenvectors satisfy $-u_1-u_2+u_3=0$), giving two independent eigenvectors $(1,0,1)$ and $(-1,1,0)$; for $r=-3$, $\vv u_3=(-1,-1,1)$. Hence
  \[
    \vv x=c_1e^{3t}\begin{bmatrix}1\\0\\1\end{bmatrix}+c_2e^{3t}\begin{bmatrix}-1\\1\\0\end{bmatrix}+c_3e^{-3t}\begin{bmatrix}-1\\-1\\1\end{bmatrix}.
  \]
\end{example*}

\noindent\textbf{When there are not enough eigenvectors (Nagle 9.5, Probs.~35--40).} The non-symmetric matrix $A=\begin{psmallmatrix}2&1\\-1&4\end{psmallmatrix}$ has $\det(A-rI)=(r-3)^2$, a double eigenvalue $3$, yet $A-3I=\begin{psmallmatrix}-1&1\\-1&1\end{psmallmatrix}$ has only a one-dimensional eigenspace $\spn\{(1,1)\}$. \ilink{thm-nagle-9-5}{Thm.~5} gives only one solution $e^{3t}(1,1)$. For a double root $r$ of a scalar equation the second solution was $te^{rt}$. The first attempt $te^{rt}\vv u_1$ fails: its derivative is $e^{rt}\vv u_1+rte^{rt}\vv u_1$, while $A(te^{rt}\vv u_1)=rte^{rt}\vv u_1$, so the term $e^{rt}\vv u_1$ is left over. An extra term $e^{rt}\vv u_2$ can absorb it, so try $\vv x=te^{rt}\vv u_1+e^{rt}\vv u_2$.

\begin{proposition*}[Second solution for a double eigenvalue (Nagle 9.5, Prob.~35)]\target{prop-double-eig-2nd}
  Let $(A-rI)\vv u_1=\vv 0$ with $\vv u_1\neq\vv 0$. Then $\vv x(t)=te^{rt}\vv u_1+e^{rt}\vv u_2$ solves $\vv x'=A\vv x$ iff $(A-rI)\vv u_2=\vv u_1$. Moreover, $e^{rt}\vv u_1$ and $\vv x(t)$ are linearly independent.
\end{proposition*}
\begin{proof}
  \begin{steps}
    \ii Since $A\vv u_1=r\vv u_1$, the product rule gives $\vv x'=e^{rt}\vv u_1+rte^{rt}\vv u_1+re^{rt}\vv u_2$, while $A\vv x=rte^{rt}\vv u_1+e^{rt}A\vv u_2$. The terms $rte^{rt}\vv u_1$ agree, and dividing the rest by $e^{rt}\neq 0$ shows $\vv x'=A\vv x \iff \vv u_1+r\vv u_2=A\vv u_2 \iff (A-rI)\vv u_2=\vv u_1$.
    \ii Since $(A-rI)(c\vv u_1)=\vv 0\neq\vv u_1$, $\vv u_2$ cannot be a multiple of $\vv u_1$; thus $W(0)=\det[\,\vv u_1\ \vv u_2\,]\neq 0$, giving independence. \qedhere
  \end{steps}
\end{proof}

\begin{proposition*}[Solvability of $(A-rI)\vv u_2=\vv u_1$ in $2\times 2$]\target{prop-u2-solvable-2x2}
  If $2\times 2$ $A$ has a double eigenvalue $r$ with one-dimensional eigenspace, then $(A-rI)\vv u_2=\vv u_1$ has a solution for every eigenvector $\vv u_1$.
\end{proposition*}
\begin{proof}
  \begin{steps}
    \ii Put $N=A-rI=\begin{psmallmatrix}a&b\\c&d\end{psmallmatrix}$. Since $r$ is a double root of $\lambda^2-(\tr A)\lambda+\det A=0$, $\tr A=2r$, so $\tr N=0$; and $\det N=0$ because $r$ is an eigenvalue. Multiplying out $N^2$:
    \[
      N^2 = \begin{bmatrix}a^2+bc&b(a+d)\\c(a+d)&d^2+bc\end{bmatrix} = (a+d)N-(ad-bc)I = (\tr N)N-(\det N)I = 0.
    \]
    \ii Since $\dim\Nul N=1$, the Rank Theorem (\ilink{fact-W1}{W1}) gives $\rank N=1$.
    \ii For every $\vv v$, $N(N\vv v)=N^2\vv v=\vv 0$, so $\Col N\subseteq\Nul N$. Both are lines, so $\Col N=\Nul N$. Hence $\vv u_1\in\Nul N=\Col N$, that is, $\vv u_1=N\vv u_2$ for some $\vv u_2$. \qedhere
  \end{steps}
\end{proof}

\begin{example*}[A double eigenvalue with one eigenvector]\target{ex-defective-2x2}
  For $A=\begin{psmallmatrix}2&1\\-1&4\end{psmallmatrix}$, $r=3$ and $\vv u_1=(1,1)$. The equation $(A-3I)\vv u_2=(1,1)$ says $-u_{21}+u_{22}=1$ (both rows); taking $\vv u_2=(0,1)$ yields
  \[
    \vv x(t) = c_1 e^{3t}\begin{bmatrix}1\\1\end{bmatrix} + c_2\left(te^{3t}\begin{bmatrix}1\\1\end{bmatrix} + e^{3t}\begin{bmatrix}0\\1\end{bmatrix}\right),
  \]
  with $W(0)=\det\begin{psmallmatrix}1&0\\1&1\end{psmallmatrix}=1\neq 0$. The vector $\vv u_2$ is called a \textbf{generalized eigenvector}. Any other choice, such as $\vv u_2=(1,2)$, differs by a multiple of $\vv u_1$ and changes only $c_1$.
\end{example*}
\begin{proposition*}[Third solution for a triple eigenvalue]\target{prop-triple-eig}
  Let $N=A-rI$ and suppose $N\vv u_1=\vv 0$, $N\vv u_2=\vv u_1$, $N\vv u_3=\vv u_2$. Then $\vv x(t)=\frac{t^2}{2}e^{rt}\vv u_1+te^{rt}\vv u_2+e^{rt}\vv u_3$ solves $\vv x'=A\vv x$.
\end{proposition*}
\begin{proof}
  \begin{steps}
    \ii Since $A=rI+N$, the equation $\vv x'=A\vv x$ is the same as $\vv x'-r\vv x=N\vv x$.
    \ii The product rule gives $\vv x'-r\vv x=te^{rt}\vv u_1+e^{rt}\vv u_2$: differentiating $\frac{t^2}{2}$, $t$ and $1$ gives $t$, $1$ and $0$, and the terms from differentiating $e^{rt}$ make up $r\vv x$.
    \ii Applying $N$ term by term gives $N\vv x=\frac{t^2}{2}e^{rt}\vv 0+te^{rt}\vv u_1+e^{rt}\vv u_2$, which is the same. \qedhere
  \end{steps}
\end{proof}
\noindent So for a triple eigenvalue with a one-dimensional eigenspace, the chain $\vv u_1,\vv u_2,\vv u_3$ gives the solutions $e^{rt}\vv u_1$, $te^{rt}\vv u_1+e^{rt}\vv u_2$, and $\frac{t^2}{2}e^{rt}\vv u_1+te^{rt}\vv u_2+e^{rt}\vv u_3$; their values at $t=0$ are $\vv u_1,\vv u_2,\vv u_3$.

\begin{howto}{How to solve $\vv x'=A\vv x$ ($A$ constant) by the eigenvalue method}\target{howto-eigenvalue-method}
  \begin{steps}
    \ii \textbf{Eigenvalues.} Solve $\det(A-rI)=0$ and record each root with its multiplicity. (For $2\times 2$: $r^2-(\tr A)r+\det A=0$.)
    \ii \textbf{Eigenvectors.} For each eigenvalue $r$, row-reduce $A-rI$ and take a basis of $\Nul(A-rI)$. Count: the number of basis vectors may be smaller than the multiplicity.
    \ii \textbf{Enough eigenvectors ($n$ in total).} The general solution is $\sum c_ie^{r_it}\vv u_i$ (\ilink{thm-nagle-9-5}{Thm.~5}). This is automatic for $n$ distinct eigenvalues (\ilink{thm-nagle-9-6}{Thm.~6}) and for symmetric $A$.
    \ii \textbf{A double eigenvalue with one eigenvector $\vv u_1$.} Solve $(A-rI)\vv u_2=\vv u_1$ and use $te^{rt}\vv u_1+e^{rt}\vv u_2$ as the second solution (\ilink{prop-double-eig-2nd}{Second solution for a double eigenvalue}).
    \ii \textbf{Initial condition at $t_0=0$.} Solve $[\,\vv u_1\ \cdots\ \vv u_n\,]\vv c=\vv x(0)$; with a generalized eigenvector the column for the second solution is $\vv u_2$, its value at $t=0$. At $t_0\neq 0$, solve $X(t_0)\vv c=\vv x_0$.
    \ii \textbf{Long-run behavior.} The term with the largest eigenvalue among those with $c_i\neq 0$ dominates as $t\to\infty$.
  \end{steps}
  \textbf{Edge cases.} An eigenvalue $0$ gives constant solutions $c\vv u$ (equilibrium points along a whole line). If $A=rI$, every nonzero vector is an eigenvector and $\vv x=e^{rt}\vv x(0)$. Complex eigenvalues give complex solutions whose real and imaginary parts are real solutions.
\end{howto}

\begin{example*}[An eigenvalue equal to $0$: a line of equilibrium points]\target{ex-zero-eig}
  For $A=\begin{psmallmatrix}1&1\\1&1\end{psmallmatrix}$, $\det(A-rI)=r^2-2r=r(r-2)$, so $r=0$ and $r=2$. $\Nul A$ is spanned by $(-1,1)$, and $A-2I$ has rows $(-1,1),(1,-1)$, giving $(1,1)$. Since $e^{0t}=1$,
  \[
    \vv x(t)=c_1\begin{bmatrix}-1\\1\end{bmatrix}+c_2e^{2t}\begin{bmatrix}1\\1\end{bmatrix}.
  \]
  Every point $c_1(-1,1)$ of the line $x_1+x_2=0$ is a constant solution. Any other solution moves away from that line along the direction $(1,1)$.
\end{example*}

\medskip\noindent\textbf{Diagonalization connection (Nagle 9.5).} With $U=[\,\vv u_1\ \cdots\ \vv u_n\,]$ ($AU=UD$, $A=UDU^{-1}$), substituting $\vv x=U\vv y$ gives $U\vv y'=AU\vv y=UD\vv y$, so $\vv y'=D\vv y$. These are $n$ separate scalar equations $y_i'=r_i y_i$, so $y_i=c_ie^{r_it}$, and $X(t)=U\operatorname{diag}(e^{r_1t},\dots,e^{r_nt})$.

\begin{shortcuts}
  \ii \textbf{IVP constants:} Solve $U\vv c=\vv x(0)$ directly ($e^{r_i\cdot 0}=1$).
  \ii \textbf{Long-run asymptotics:} The largest real eigenvalue $r_{\max}$ (with $c_i\neq 0$) dominates as $t\to\infty$; all $r_i<0\implies \vv x(t)\to\vv 0$.
  \ii \textbf{$2\times 2$ shortcut:} Characteristic equation $r^2-(\tr A)r+\det A=0$; a nonzero row $(a,b)$ of $A-rI$ gives the eigenvector $(b,-a)$.
  \ii \textbf{Triangular $A$:} Eigenvalues sit on the diagonal (\ilink{fact-W1}{W1}).
  \ii \textbf{Sum/product check:} Eigenvalues always sum to $\tr A$ and multiply to $\det A$.
  \ii \textbf{Defective $2\times 2$:} Count $\dim\Nul(A-rI)$ first; if $1$, use $te^{rt}\vv u_1+e^{rt}\vv u_2$ with $(A-rI)\vv u_2=\vv u_1$.
\end{shortcuts}

\begin{trapbox}
  \ii \textbf{A repeated eigenvalue may or may not have enough eigenvectors:} the symmetric example has two for $r=3$, the matrix $\begin{psmallmatrix}2&1\\-1&4\end{psmallmatrix}$ has one. Always compute $\dim\Nul(A-rI)$.
  \ii \textbf{Never drop $e^{rt}\vv u_2$ on a defective eigenvalue:} $te^{rt}\vv u_1$ alone satisfies $\vv x'-A\vv x=e^{rt}\vv u_1\neq\vv 0$.
  \ii \textbf{Triple eigenvalue with geometric multiplicity~$2$:} $(A-rI)\vv w=\vv u$ works \emph{only} for certain eigenvectors $\vv u$ in $\Col(A-rI)\cap\Nul(A-rI)$. For the matrix
  \[
    A=\begin{bmatrix}1&1&0\\0&1&0\\0&0&1\end{bmatrix}
  \]
  with triple eigenvalue $1$ and eigenvectors $\vv e_1,\vv e_3$, $(A-I)\vv w=(w_2,0,0)^{\mathsf T}$, so $(A-I)\vv w=\vv e_3$ has no solution while $\vv u=\vv e_1,\ \vv w=\vv e_2$ works.
  \ii \textbf{Match $\vv x(0)$ using the full combination $\sum c_i\vv u_i$:} Never assume $\vv x(0)$ is an eigenvector.
  \ii \textbf{Characteristic equation is $\det(A-rI)=0$, \emph{not} $\det(A)-r=0$.}
  \ii \textbf{An eigenvalue $0$ is allowed:} its eigenvectors give constant solutions, not the zero solution.
\end{trapbox}

\section{Proofs of the facts used}\target{sec-proofs}

\subsection{Existence and uniqueness for linear systems}\target{proof-existence}
This subsection proves the \ilink{thm-nagle-9-2}{existence and uniqueness theorem}: if $A(t)$ and $\vv f(t)$ are continuous on an open interval $I$ containing $t_0$, then $\vv x'=A(t)\vv x+\vv f(t)$, $\vv x(t_0)=\vv x_0$ has exactly one solution on $I$.

\begin{ideabox}
  Integrating the equation from $t_0$ to $t$ turns it into $\vv x(t)=\vv x_0+\int_{t_0}^t\bigl(A\vv x+\vv f\bigr)ds$, a statement that involves no derivative of $\vv x$. Its right side can be computed for any guess $\vv x$, so start from the constant guess $\vv x_0$, substitute it into the right side to get a better guess, and repeat. Each round improves the error by a factor that behaves like $K|t-t_0|/k$ at round $k$; these factors multiply to $(K|t-t_0|)^k/k!$, the terms of the series for $e^{K|t-t_0|}$, so the guesses converge. The same estimate applied to the difference of two solutions shows that the difference is $0$.
\end{ideabox}

\noindent\textbf{Size of a vector and of a matrix.} For $\vv v\in\RR^n$ let $|\vv v|=\max_i|v_i|$, the largest absolute value of an entry. For an $n\times n$ matrix $B$ let $\|B\|=\max_i\sum_j|b_{ij}|$, the largest absolute row sum. Then:
\begin{itemize}
  \ii $|\vv u+\vv v|\le|\vv u|+|\vv v|$, because $|u_i+v_i|\le|u_i|+|v_i|\le|\vv u|+|\vv v|$ for every $i$.
  \ii $|B\vv v|\le\|B\|\,|\vv v|$, because for every $i$
  \[
    |(B\vv v)_i|=\Bigl|\sum_j b_{ij}v_j\Bigr|\le\sum_j|b_{ij}|\,|\vv v|\le\|B\|\,|\vv v|.
  \]
  \ii The size of an integral is at most the integral of the size:
  \[
    \Bigl|\int_a^b\vv g(s)\,ds\Bigr|\le\Bigl|\int_a^b|\vv g(s)|\,ds\Bigr|,
  \]
  because each entry satisfies $|\int g_i|\le|\int|g_i||$ (\ilink{fact-W5}{W5}) and $|g_i(s)|\le|\vv g(s)|$.
\end{itemize}

\begin{proof}[Proof of existence and uniqueness]
  \begin{steps}
    \ii \textbf{The integral form.} A continuous $\vv x$ solves the IVP on an interval $J\ni t_0$ if and only if
    \[
      \vv x(t)=\vv x_0+\int_{t_0}^t\bigl(A(s)\vv x(s)+\vv f(s)\bigr)\,ds\qquad\text{for all }t\in J.
    \]
    \textbf{Assume $\vv x$ solves the IVP.} Integrating $\vv x'=A\vv x+\vv f$ from $t_0$ to $t$ and using $\vv x(t_0)=\vv x_0$ gives the formula. \textbf{Assume the formula holds.} The integrand is continuous, so by the fundamental theorem of calculus (\ilink{fact-W5}{W5}) the right side is differentiable with derivative $A\vv x+\vv f$, and at $t=t_0$ it equals $\vv x_0$.
    \ii \textbf{Constants on a closed interval.} Fix a closed bounded interval $J=[\alpha,\beta]\subseteq I$ containing $t_0$, and let $\ell=\beta-\alpha$. Each entry of $A$ is continuous on $J$, so it is bounded there (\ilink{fact-W5}{W5}); hence $K=\max_{t\in J}\|A(t)\|$ is finite. For $t,s\in J$, $|t-t_0|\le\ell$.
    \ii \textbf{Uniqueness on $J$.} Let $\vv x,\vv y$ both solve the IVP on $J$ and put $\vv z=\vv x-\vv y$. Subtracting the integral forms, $\vv z(t)=\int_{t_0}^tA(s)\vv z(s)\,ds$. Let $m=\max_J|\vv z|$ (finite, since $\vv z$ is continuous). We show by induction on $k$ that
    \[
      |\vv z(t)|\le m\,\frac{K^k|t-t_0|^k}{k!}\qquad(t\in J).
    \]
    For $k=0$ this is the definition of $m$. If it holds for $k$, then
    \[
      |\vv z(t)|\le\Bigl|\int_{t_0}^tK\,m\,\frac{K^k|s-t_0|^k}{k!}\,ds\Bigr|=m\,\frac{K^{k+1}|t-t_0|^{k+1}}{(k+1)!}.
    \]
    So $|\vv z(t)|\le m(K\ell)^k/k!$ for every $k$. These numbers are terms of the convergent series for $e^{K\ell}$, so they tend to $0$ (\ilink{fact-W5}{W5}), and $\vv z(t)=\vv 0$.
    \ii \textbf{Successive approximations.} Define $\vv x_0(t)=\vv x_0$ and
    \[
      \vv x_{k+1}(t)=\vv x_0+\int_{t_0}^t\bigl(A(s)\vv x_k(s)+\vv f(s)\bigr)\,ds\qquad(k=0,1,2,\dots).
    \]
    Each $\vv x_k$ is continuous on $J$ (an integral of a continuous function is differentiable). Let $M=\max_J|\vv x_1-\vv x_0|$. Since $\vv x_{k+1}-\vv x_k=\int_{t_0}^tA(\vv x_k-\vv x_{k-1})\,ds$ has the same shape as the equation for $\vv z$ in step~3, the same induction gives
    \[
      |\vv x_{k+1}(t)-\vv x_k(t)|\le M\,\frac{K^k|t-t_0|^k}{k!}\le M\,\frac{(K\ell)^k}{k!}\qquad(t\in J).
    \]
    \ii \textbf{The approximations converge, at the same rate for every $t$.} For each $t$, every entry of $\vv x_N(t)=\vv x_0+\sum_{k<N}(\vv x_{k+1}(t)-\vv x_k(t))$ is a partial sum of a series whose terms are bounded in absolute value by $M(K\ell)^k/k!$, so the series converges absolutely by comparison. Call the limit $\vv x(t)$. Its distance from the $N$th approximation is bounded by the tail:
    \[
      |\vv x(t)-\vv x_N(t)|\le\varepsilon_N=M\sum_{k\ge N}\frac{(K\ell)^k}{k!},
    \]
    and $\varepsilon_N\to 0$ as $N\to\infty$ (\ilink{fact-W5}{W5}). The bound $\varepsilon_N$ does not depend on $t$.
    \ii \textbf{The limit is continuous.} For $t,s\in J$,
    \[
      |\vv x(t)-\vv x(s)|\le|\vv x(t)-\vv x_N(t)|+|\vv x_N(t)-\vv x_N(s)|+|\vv x_N(s)-\vv x(s)|\le 2\varepsilon_N+|\vv x_N(t)-\vv x_N(s)|.
    \]
    Given $\varepsilon>0$, choose $N$ with $2\varepsilon_N<\varepsilon/2$, then use the continuity of $\vv x_N$ to make the last term less than $\varepsilon/2$ for $s$ close to $t$.
    \ii \textbf{The limit satisfies the integral form.} By the bounds above,
    \[
      \Bigl|\vv x_0+\int_{t_0}^t(A\vv x+\vv f)\,ds-\vv x_{N+1}(t)\Bigr|=\Bigl|\int_{t_0}^tA(\vv x-\vv x_N)\,ds\Bigr|\le K\ell\,\varepsilon_N\to 0,
    \]
    while $\vv x_{N+1}(t)\to\vv x(t)$. Hence $\vv x(t)=\vv x_0+\int_{t_0}^t(A\vv x+\vv f)\,ds$, and by step~1, $\vv x$ solves the IVP on $J$.
    \ii \textbf{From closed intervals to all of $I$.} Choose closed bounded intervals $J_1\subseteq J_2\subseteq\cdots$ in $I$, each containing $t_0$, whose union is $I$ (for $I=(p,q)$ with $p<t_0<q$, take $J_k=[p+\delta_k,\,q-\delta_k]$ with $\delta_k\downarrow 0$, or $[t_0-k,t_0+k]$ for an infinite end). Let $\vv x^{(k)}$ be the solution on $J_k$ from steps~4--7. By step~3 on $J_k$, $\vv x^{(k+1)}=\vv x^{(k)}$ on $J_k$. So setting $\vv x(t)=\vv x^{(k)}(t)$ for any $k$ with $t\in J_k$ defines one function on $I$, and it solves the IVP near every point of $I$. If $\vv y$ is another solution on $I$, step~3 on each $J_k$ gives $\vv y=\vv x$ on $J_k$, hence on $I$. \qedhere
  \end{steps}
\end{proof}

\begin{example*}[Successive approximations for $x'=x$, $x(0)=1$]\target{ex-picard}
  Here $n=1$, $A=[1]$, $f=0$. The approximations are $x_0=1$, $x_1=1+\int_0^t1\,ds=1+t$, $x_2=1+\int_0^t(1+s)\,ds$, which is $1+t+\frac{t^2}{2}$, and in general $x_k=1+t+\dots+\frac{t^k}{k!}$, the partial sums of the series of $e^t$. Their limit $e^t$ is the solution.
\end{example*}

\subsection{Abel's formula for $n\times n$ systems}\target{proof-abel-n}
This subsection proves \ilink{prop-abel-formula}{Abel's formula} for $n\times n$ systems: if the columns of $X(t)$ solve $\vv x'=A(t)\vv x$, then $W=\det X$ satisfies $W'=(\tr A)W$.

\begin{ideabox}
  The determinant is a sum of products with one entry from each row, so by the product rule its derivative is a sum of $n$ determinants, each with one row differentiated. The equation $X'=AX$ says that a differentiated row is a combination of the rows of $X$; in a determinant, the other rows of that combination can be removed by row replacement, which leaves only the diagonal coefficient $a_{ii}$.
\end{ideabox}

\begin{proof}[Proof of Abel's formula for $n\times n$]
  \begin{steps}
    \ii \textbf{The determinant as a sum of products.} Expanding $\det X$ by cofactors along row~1, then each minor along its first row, and so on (\xlink{math54}{2026-09-14-matrix-algebra-inverses}{thm-lay-3-1}{cofactor expansion, Week~4}), writes $\det X$ as a sum of terms $\pm x_{1j_1}x_{2j_2}\cdots x_{nj_n}$, each a product of one entry from each row, with signs that do not depend on $t$.
    \ii \textbf{Differentiating.} By the product rule, the derivative of $x_{1j_1}\cdots x_{nj_n}$ is the sum over $i$ of the same product with the factor from row $i$ replaced by its derivative. Collecting, for each $i$, the terms in which row $i$ was differentiated gives $W'=\sum_{i=1}^n\det X_i$, where $X_i$ is $X$ with row $i$ replaced by row $i$ of $X'$.
    \ii \textbf{Row $i$ of $X'$.} Each column satisfies $\vv x_j'=A\vv x_j$, so $X'=AX$, and row $i$ of $AX$ is $\sum_{k}a_{ik}(\text{row }k\text{ of }X)$.
    \ii \textbf{Removing the other rows.} In $X_i$, subtract $a_{ik}$ times row $k$ from row $i$ for every $k\neq i$. Row replacements do not change the determinant (\xlink{math54}{2026-09-21-determinants-vector-spaces}{thm-lay-3-3}{row operations and the determinant, Week~5}), and they leave $a_{ii}(\text{row }i\text{ of }X)$ in row $i$. Factoring $a_{ii}$ out of that row gives $\det X_i=a_{ii}\det X$.
    \ii \textbf{Conclusion.} $W'=\sum_ia_{ii}W=(\tr A)W$, and step~2 of the \ilink{prop-abel-formula}{$2\times 2$ proof} turns this into $W(t)=W(t_0)\exp\int_{t_0}^t\tr A(s)\,ds$. \qedhere
  \end{steps}
\end{proof}

\section{Homework this week}
\begin{itemize}
  \ii \textbf{HW12, Nagle Section 9.1:} 2, 4, 7, 12 (normal form, Section~2.1).
  \ii \textbf{Nagle Section 9.4:} 1, 3, 7, 10, 14, 16, 21, 25, 31 (Wronskians \& fundamental matrices, Sections~2.1--2.2).
  \ii \textbf{Nagle Section 9.5:} 12, 13, 17(a,b,c), 20, 25, 31, 35 (eigenvalue method, Section~2.3).
  \ii \textbf{HW12 also includes Nagle Section 4.5:} 2(b), 8, 10, 11, 18, 29, 35 (\xlink{math54}{2026-11-09-undetermined-coefficients-superposition}{thm-nagle-4-4}{Week~12}, Section~2.2).
\end{itemize}

\section{Problem set}
\begin{prob}\target{prob-P1}
  \textbf{P1.}\ Write $y'''-2y''+y'-3y=e^t$ as a first-order system in normal form $\vv x'=A\vv x+\vv f$.
\end{prob}
\begin{prob}\target{prob-P2}
  \textbf{P2.}\ Write the system $x''+2x-y=0,\ y''+3y'-x=0$ in normal form.
\end{prob}
\begin{prob}\target{prob-P3}
  \textbf{P3.}\ Show that $e^{-2t}(1,-2)$ and $e^{-3t}(1,-3)$ solve the system with matrix $A$ below, form a fundamental matrix, and give the general solution. Which scalar equation does this system come from?
  \[
    A=\begin{bmatrix}0&1\\-6&-5\end{bmatrix}.
  \]
\end{prob}
\begin{prob}\target{prob-P4}
  \textbf{P4.}\ Are $(1,1+t)$ and $(t,t+t^2)$ linearly independent on $(-\infty,\infty)$? Compute their Wronskian. Can they both be solutions of one system $\vv x'=A(t)\vv x$ with $A$ continuous on $(-\infty,\infty)$?
\end{prob}
\begin{prob}\target{prob-P5}
  \textbf{P5.}\ Find the eigenvalues and eigenvectors of the matrix below, the general solution of $\vv x'=A\vv x$, and the solution with $\vv x(0)=(1,0)$.
  \[
    A=\begin{bmatrix}4&-2\\1&1\end{bmatrix}.
  \]
\end{prob}
\begin{prob}\target{prob-P6}
  \textbf{P6.}\ Find the general solution of $\vv x'=A\vv x$ for the symmetric matrix
  \[
    A=\begin{bmatrix}1&0&2\\0&3&0\\2&0&1\end{bmatrix}.
  \]
\end{prob}
\begin{prob}\target{prob-P7}
  \textbf{P7.}\ The matrix below has a double eigenvalue with only one independent eigenvector. Find two independent solutions of $\vv x'=A\vv x$.
  \[
    A=\begin{bmatrix}3&1\\-1&1\end{bmatrix}.
  \]
\end{prob}
\begin{prob}\target{prob-P8}
  \textbf{P8.}\ For $\vv x'=A\vv x$ with the matrix below, describe the long-run behavior of all solutions, and find the solutions that tend to $\vv 0$.
  \[
    A=\begin{bmatrix}1&2\\2&1\end{bmatrix}.
  \]
\end{prob}
\end{document}
